% ==================
% Zielsetzung
% ==================

\section{Zielsetzung}

Das primäre Ziel dieser Diplomarbeit ist die Realisierung von „LearnHub“, einer spezialisierten Web-Plattform zur gezielten Matura-Vorbereitung an der HTL Leonding. Das Gesamtprojekt „Maturatrainer“ ist dabei in zwei komplementäre Arbeitsbereiche unterteilt: Während das Partnerteam das Content-Management mit der Erstellung von Fragenkatalogen, Themenpools und dem Material-Upload verantwortet, konzentriert sich die vorliegende Arbeit auf die methodische Nutzbarmachung dieser Inhalte. Die technische Umsetzung erfolgt als moderne Web-Applikation mittels \textbf{Angular} im Frontend, einer \textbf{Quarkus}-REST-API im Backend sowie einer \textbf{MySQL}-Datenbank.

Im Fokus der Entwicklung stehen funktionale Schwerpunkte, die im praktischen Teil dieser Arbeit vertieft behandelt werden. Das System bietet einen \textbf{intelligenten Übungsmodus}, der auf dem \textit{Leitner-System} sowie \textit{Spaced Repetition} basiert, um einen langfristigen Wissensaufbau zu garantieren. Ergänzend dazu ermöglicht eine \textbf{konfigurierbare Prüfungssimulation} mit variablen Zeitlimits und spezieller Navigationslogik eine realitätsnahe Abbildung der Matura-Situation. 

Um den Lernfortschritt objektiv bewerten zu können, bereitet die Anwendung \textbf{dynamische Lernstatistiken} in visuellen Dashboards auf, wodurch spezifische Wissenslücken sofort identifizierbar werden. Eine \textbf{individuelle Lernsteuerung} erlaubt es den Nutzer:innen zudem, gezielt aus den vom Partnerteam bereitgestellten Inhalten zu wählen. Abgerundet wird das System durch \textbf{Gamification-Elemente} wie „Streaks“ und kontinuierliche Feedback-Zyklen, die die Eigenverantwortung und eine regelmäßige Nutzung der Plattform fördern.